% ECONOMICS_PROBLEM: {"id": "H22-444", "title": "Corporate-tax incidence under global profit shifting", "jel_code": "H22", "source_id": "OP-0043"}
% !TeX program = xelatex
\documentclass[11pt,a4paper]{article}
\usepackage[margin=23mm]{geometry}
\usepackage{fontspec}
\setmainfont{Latin Modern Roman}
\usepackage{amsmath,amssymb,mathtools}
\usepackage{xcolor,enumitem,booktabs,longtable,array,xurl,hyperref}
\definecolor{muted}{HTML}{53636C}
\hypersetup{colorlinks=true,urlcolor=blue,linkcolor=blue}
\setlength{\parindent}{0pt}
\setlength{\parskip}{5pt}
\newcommand{\R}{\mathbb R}
\newcommand{\E}{\mathbb E}
\newcommand{\N}{\mathbb N}
\newcommand{\ind}{\mathbf 1}
\newcommand{\Var}{\operatorname{Var}}
\newcommand{\Cov}{\operatorname{Cov}}
\newcommand{\supp}{\operatorname{supp}}
\DeclareMathOperator*{\argmin}{arg\,min}
\DeclareMathOperator*{\argmax}{arg\,max}
\newcommand{\entrymeta}[3]{{\sffamily\small\color{muted}\textbf{\texttt{#1}}\quad|\quad #2\par #3\par}}
\newcommand{\Problemref}[1]{\texttt{#1}}
\newcommand{\JELref}[2]{\texttt{#2} (JEL \texttt{#1})}
\begin{document}
\section*{Corporate-tax incidence under global profit shifting}
\textbf{Problem ID:} \texttt{H22-444}\quad \textbf{JEL:} \texttt{H22}\par
% BEGIN ECONOMICS STATEMENT
\label{problem:OP-0043}
\label{atlas:prob:P3}
\noindent\textbf{Original atlas identifier:} P3; entry 43. \textbf{Field:} Public finance and taxation. \textbf{Original tractability label:} Hard.

\noindent\textbf{Classification:} Parameterized theoretical/empirical research agenda; not a certified unsolved theorem.

\subsection*{Definitions and assumptions}
Consider countries $j=1,\ldots,J$ linked by trade, capital ownership, and multinational production. A policy vector $p$ includes corporate rates, apportionment, transfer-pricing enforcement, and any explicitly dated minimum-tax constraints. Each structural model $m$ determines real activity, booked profits, wages, consumer prices, land rents, and ownership returns; these are measured separately. Fix equilibrium selection $s_m$ and a population partition into disjoint beneficiary groups $g=1,\ldots,G$, allowing a household's several income sources to be combined before group assignment. Let $EV_g(p;m)$ be aggregate money-metric equivalent variation for group $g$, measured against policy $p_0$, and let $R_j$ and $C_j$ be public revenue and administrative resource cost. Select a differentiable feasible path $p(\epsilon)$ with $p(0)=p_0$. Assume also that every composite $\epsilon\mapsto EV_g(p(\epsilon);m)$ is differentiable at zero when the marginal estimand below is used; differentiability of the policy path alone does not ensure this.

\subsection*{Formal research question}
\emph{Editorial formulation: the mathematical target below is not a theorem quoted from the references.}

Define marginal private incidence by
\[
 I_g(m)=-\left.\frac{d}{d\epsilon}EV_g(p(\epsilon);m)\right|_{\epsilon=0}.
\]
At kinks or equilibrium jumps, use a specified one-sided derivative only when its finite limit exists, or the discrete contrast $-[EV_g(p_1;m)-EV_g(p_0;m)]$ for two stated feasible policies $p_0,p_1$. Identify the vector $(I_g(m))_{g=1}^{G}$, its horizon dependence, and the joint response of each $R_j-C_j$. When $\sum_g I_g(m)\neq0$, normalized shares $I_g/\sum_h I_h$ may be reported, retaining signs; they need not lie in $[0,1]$. Compare global and domestic welfare changes using explicitly chosen group weights $w_g$ and public-fund values $\lambda_j$: $\sum_gw_g\Delta EV_g+\sum_j\lambda_j\Delta(R_j-C_j)$. Establish what matched firm--worker--ownership data and exogenous tax variation identify independently of restrictions on profit shifting. If multiple models fit the data, deliver incidence bounds rather than a single bearer.

\subsection*{Interpretation and scope}
A tax dollar is a transfer to government, whereas lost production and administrative effort are resource costs; fiscal and private incidence should not be conflated. Short-run rents, long-run investment, and international relocation can generate different incidence vectors. Legal booking jurisdiction is not a welfare beneficiary and does not reveal whose wages or returns change. Local tax reforms identify local responses under their assumptions, not automatically worldwide coordination counterfactuals. Consumption-price effects and pension ownership must be incorporated when the group called “workers” also owns capital. Minimum-tax rules must be described by the legal version actually studied, with a separately checked effective date. The research extension is international joint identification and equilibrium assessment, not a presumption that classical or local incidence models have no solutions.

\subsection*{Source audit and status}
All three original academic references are identifiable. [S1] is a classical benchmark; [S2] and [S3] are setting-specific empirical/structural results. The publisher of [S2] links a later comment and reply, which should accompany replication. These sources do not establish contemporary international incidence shares.

\subsection*{References}
\begin{enumerate}[label={[S\arabic*]},leftmargin=*,itemsep=0.6em]
\item Arnold C. Harberger (1962). The Incidence of the Corporation Income Tax. Journal of Political Economy 70(3), 215--240. DOI: 10.1086/258636.
\url{https://www.journals.uchicago.edu/doi/10.1086/258636}
\newline\emph{Locator and support limits:} Article, pp. 215--240: closed-economy general-equilibrium incidence benchmark; not a multinational shifting model.
\item Juan Carlos Suárez Serrato and Owen Zidar (2016). Who Benefits from State Corporate Tax Cuts? A Local Labor Markets Approach with Heterogeneous Firms. American Economic Review 106(9), 2582--2624. DOI: 10.1257/aer.20141702.
\url{https://www.aeaweb.org/articles?id=10.1257/aer.20141702}
\newline\emph{Locator and support limits:} Abstract and article: local workers, landowners, and firm owners; publisher links a 2023 comment and reply, so estimates merit that context.
\item Clemens Fuest, Andreas Peichl, and Sebastian Siegloch (2018). Do Higher Corporate Taxes Reduce Wages? Micro Evidence from Germany. American Economic Review 108(2), 393--418. DOI: 10.1257/aer.20130570.
\url{https://pubs.aeaweb.org/doi/abs/10.1257/aer.20130570}
\newline\emph{Locator and support limits:} Abstract and German municipal analysis: wage incidence; does not estimate a global minimum-tax equilibrium.
\end{enumerate}

\subsection*{Common conventions: Source mathematical conventions}

Unless an entry states otherwise, agent and firm sets are finite. State and action spaces are measurable, strategies and outcome maps are measurable, probability laws are specified, and expectations exist for the targets under discussion. Infinite-horizon discount factors lie in $(0,1)$ and discounted objectives are finite. Norms, tolerances and complexity input sizes must be specified for computational comparisons. Constants or primitives that appear locally are inputs, not empirical facts asserted by this catalogue.

\textbf{Identification.} A statistical model $m\in\mathcal M$ implies an observable law $P_m$ and target $\theta(m)$. At law $P$, its identified set is
\[
\Theta(P;\mathcal M)=\{\theta(m):m\in\mathcal M,\ P_m=P\}.
\]
Point identification holds when this set is a singleton, or equivalently when $P_{m_1}=P_{m_2}$ implies $\theta(m_1)=\theta(m_2)$ for all compatible models. Sharp bounds mean bounds attained, or approached when the boundary is not attained, by compatible models. Writing this set is a definition, not a solution: the substantive task is to establish informative restrictions and derive or estimate the set. An unrestricted-confounding model may give only trivial bounds.

\textbf{Inference.} Identification, estimability and valid uncertainty quantification are separate questions. Fix $\alpha\in(0,1)$. A confidence procedure $C_n$ over a declared law class $\mathcal P$ has asymptotic uniform coverage at level $1-\alpha$ when
\[
\liminf_{n\to\infty}\inf_{P\in\mathcal P}\Pr_P\{\theta(P)\in C_n\}\geq1-\alpha.
\]
For set-identified targets the coverage event must be defined separately, for example coverage of the full identified set. If an entry uses identification-indexed models instead of uniquely indexed laws, the same coverage convention is applied over those models. Pointwise limits do not establish uniform validity, and asymptotic validity does not guarantee finite-sample precision. Sample sizes, dependence structures and weak-identification sequences are part of each application.

\textbf{Counterfactuals.} $Y_i(a)$ is a potential outcome under a specified intervention, including its timing, implementation and versions. It is not $Y_i$ conditional on voluntarily choosing $a$. A full intervention vector permits interference; shorthand scalar treatment notation does not silently impose no interference. Assignment, selection, attrition, anticipation and measurement must be specified. A claim about an instrument's local effect is not automatically a population policy effect.

\textbf{Economic equilibrium.} A counterfactual model includes best responses, constraints, feasibility, market clearing, state transitions and expectations consistent with the stated information. When several equilibria exist, either a declared selector is an input or the target is set-valued. Comparisons across policy regimes must use the stated selection convention. Reduced-form relationships are not presumed invariant after an equilibrium-changing intervention.

\textbf{Optimization and welfare.} Optimization is over the stated feasible set. If attainment is not established, $\sup$ or $\inf$ and $\varepsilon$-optimal choices replace an asserted maximizer or minimizer. For example, with value $V=\sup_{a\in\mathcal A}W(a)<\infty$, an $\varepsilon$-optimal set is $\{a\in\mathcal A:W(a)\geq V-\varepsilon\}$ for $\varepsilon>0$. Benefits, costs, risk and health or environmental outcomes can be aggregated only using declared units and conversion weights. Interpersonal utility weights, the treatment of fiscal transfers and foreign profits, discounting, and the behavioral welfare criterion are explicit inputs. A comparative-static derivative additionally requires existence and appropriate differentiability of the selected counterfactual.

\textbf{Sources.} Local labels [S1], [S2], and so forth restart in every question. Locators identify the relevant abstract, model, section or result at the level actually checked; a bibliographic or abstract-level verification is not represented as inspection of an entire proof. Working papers are distinguished from later publications and changing versions. Source summaries are paraphrased. All URL checks and editorial formulations refer to the audit date stated above.
% END ECONOMICS STATEMENT
\end{document}
